% META: {"id": "1NSI-TAB-EX-012", "chapitre": "1NSI-TABLES", "type_objet": "exercice", "capacites": ["1NSI-TABLES-C2"], "parcours": 2, "competences": ["programmer"], "duree_min": 8, "mode_creation": "ex_nihilo", "sources_inspiration": [], "fichier_tex": "chapitres/1NSI-TABLES/exercices/1NSI-TAB-EX-012.tex", "corrige_tex": "chapitres/1NSI-TABLES/corriges/1NSI-TAB-CO-012.tex", "status": "approved", "capacites_codes": ["C2"], "methodes": ["M1"]}
% BEGIN-VERIFY
% notes = [
%     {"nom": "Amine", "matiere": "Maths", "note": 14},
%     {"nom": "Amine", "matiere": "NSI", "note": 16},
%     {"nom": "Lina", "matiere": "Maths", "note": 12},
%     {"nom": "Lina", "matiere": "NSI", "note": 18},
%     {"nom": "Omar", "matiere": "Maths", "note": 9},
%     {"nom": "Omar", "matiere": "NSI", "note": 11},
% ]
% resultat = [ligne for ligne in notes if ligne["matiere"] == "NSI" and ligne["note"] >= 15]
% assert [ligne["nom"] for ligne in resultat] == ["Amine", "Lina"]
% END-VERIFY
\begin{exercice}{1NSI-TAB-EX-002}{1}{12}
On dispose de la table \lstinline{notes} d'un contrôle continu :
\begin{python}
notes = [
    {"nom": "Amine", "matiere": "Maths", "note": 14},
    {"nom": "Amine", "matiere": "NSI", "note": 16},
    {"nom": "Lina", "matiere": "Maths", "note": 12},
    {"nom": "Lina", "matiere": "NSI", "note": 18},
    {"nom": "Omar", "matiere": "Maths", "note": 9},
    {"nom": "Omar", "matiere": "NSI", "note": 11},
]
\end{python}
\begin{enumerate}
  \item Écrire une compréhension de liste renvoyant les lignes où la matière est \lstinline{"NSI"} et la note est supérieure ou égale à $15$.
  \item Combien de lignes obtient-on ? Donner les noms correspondants.
  \item Écrire une compréhension de liste renvoyant les noms des élèves ayant une note strictement inférieure à $10$ (peu importe la matière).
\end{enumerate}
\end{exercice}
